\SourceLeaf{245}%
\Needspace{6\baselineskip}\begin{center}\emph{[\textup{\textsc{Gerling}} an \textup{\textsc{Gauss}}. Marburg, 26. Februar 1844.]}\end{center}

\{\mbox{.\ .\ .\ .\ .\ .\ .} Die Angaben zum Buch von \textsc{Bolyai} habe ich bereits genutzt, um es zu bestellen. Mit dem ausführlichen Titel konnte sich unsere Bibliothek unmittelbar an die Buchhändler wenden. Zugleich haben wir das Buch von \textsc{Lobatschewski} aus Berlin bestellt, sodass wir die Bücher hoffentlich bald erhalten. Ich danke Ihnen nochmals herzlich für Ihre freundliche Mühe. Die russischen Abhandlungen sind mir freilich unzugänglich, da ich selbst das Buchstabieren, das ich als Student gelernt hatte, wieder vergessen habe. Den Aufsatz in \textsc{Crelle}, Band 17, habe ich inzwischen auch nachgesehen. Darin hängt allerdings alles an dem j'ai démontré. --- Das russische Steppenland scheint also doch ein geeigneter Boden für diese Spekulationen zu sein, denn \textsc{Schweikart} (jetzt in Königsberg) ersann seine »Astral-Geometrie«, während er in Charkow war. \mbox{.\ .\ .\ .\ .\ .\ .}\}

\par\bigskip\Needspace{6\baselineskip}

\begin{center}\textsc{Gauss} an \textsc{Schumacher}. Göttingen, 28. November 1846.\end{center}

\mbox{.\ .\ .\ .\ .\ .\ .\ .\ .} Ich hatte kürzlich Anlass, die kleine Schrift von \textsc{Lobatschewski} (Geometrische Untersuchungen zur Theorie der Parallellinien. Berlin 1840, bei G. Fincke. 4 Druckbogen) erneut durchzusehen. Sie enthält die Grundzüge derjenigen Geometrie, die gelten müsste und streng folgerichtig gelten könnte, wenn die euklidische Geometrie nicht die wahre ist. Ein gewisser \textsc{Schweikart}\HistoricalNote{*)}{*)}{Früher in Marburg, jetzt Professor der Rechtswissenschaft in Königsberg.} nannte eine solche Geometrie Astralgeometrie, \textsc{Lobatschewski} imaginäre Geometrie. Sie wissen, dass ich schon seit 54 Jahren (seit 1792) dieselbe Überzeugung habe (mit einer gewissen späteren Erweiterung, die ich hier nicht erwähnen will). Für mich inhaltlich Neues habe ich
